% Folder 119, batch 14 — archivists' pages 261–280.
% Pass: Opus 5.5 (claude-opus-5-5), 2026-10-03 — first pass, unchecked against the pages by a human.
%
% Contents of the batch, item by item:
% - pp. 261–270: end of the retyped and expanded typescript « Le maître-enseignant
%   et le maître-chercheur… » (begun at p. 254, batch 13), his own pagination 8–17.
%   Undated on these pages (context: spring 1968). Grothendieck's text, typed.
%   TRANSCRIBED in full. Typed underlining is rendered \emph{}. A few ink retouches
%   on the copy (pp. 264, 268) are of undetermined hand and are only noted.
% - p. 271: blank archivists' divider — no \page{}.
% - pp. 272–273: Orsay commission circular on the two doctorates (collective
%   stencil, undated; next meeting « le mardi 11 juin »). Third-party / collective
%   text: NOT transcribed, note only. The pencil notes at the foot of p. 273 are
%   in a rapid cursive judged compatible with his hand (attribution probable, not
%   certain) and echo his own typescript's themes: TRANSCRIBED, flagged.
% - p. 274: blank archivists' divider — no \page{}.
% - pp. 275–280: « Rapport de la Commission relative au recrutement des
%   enseignants » (Orsay; manuscript stencil in another hand, pp. 275–279; p. 280
%   a typed p. -6-), undated (refers to « la rentrée d'octobre 1968 »).
%   Collective and third-party (RIGHTS.md): NOT transcribed, \page{275} plus one
%   note with a three-line summary.
% The survey map supplied for the batch agrees with the pages.
\documentclass[11pt,a4paper]{article}
\input{../preamble/grothendieck}

\folder{119}
\batch{14}
\pages{261}{280}
\foldertitle{Esquisse d’un programme [dossier constitué en vue d'une candidature au CNRS (1984)] : copies de tapuscrits et de tapuscrits annotés (1972-1974, 1978-1979, 1984, 1990-1991, s.d.), tirés à part (1971), notes et copies de note manuscrites (1986, s.d.), lettres (1968, 1972, 1991).}
\dating{1968-1991}
\watermark{Édition de démonstration}

\begin{document}

\section{Le maître-enseignant et le maître-chercheur (fin)}

\page{261}
\note{Fin du tapuscrit commencé p. 254 (lot 13), pages 8 à 17 de sa propre pagination. La page s'ouvre au milieu d'une phrase venue de la p. 7 du tapuscrit (« la tête bien faite »). Les passages soulignés à la machine sont rendus en italique.}

bien faite" par opposition à la "tête bien pleine". Cela ne suppose pas nécessairement que le maître-enseignant soit au courant, à chaque moment, de tous les plus récents progrès de sa spécialité, ce qui demanderait de sa part un effort de mise au courant incompatible avec ses tâches d'enseignement. Mais les longues vacances universitaires lui permettent tout au moins de se mettre au courant périodiquement des plus importants parmi ces progrès, ou de contribuer, dans la mesure de ses moyens, à l'avancement de sa science. Ces moyens personnels sont évidemment très variables de l'un à l'autre. Au point de vue du maître-enseignant où nous nous plaçons ici, l'étendue des dons pour la recherche devient d'ailleurs un point secondaire, dans la mesure où sa participation à l'édification de la science n'est pas conçue comme une fin en soi, mais comme un moyen parmi d'autres pour garder vivant l'enseignement,- étant entendu que cela suppose d'abord que l'enseignant lui-même reste intellectuellement vivant, c'est-à-dire intellectuellement actif.

Dans tout cela, la vocation principale reste l'enseignement, et en définitive \emph{le sujet principal d'intérêt du maître-enseignant est l'étudiant lui-même}. Il est clair qu'une telle vocation ne pourra s'épanouir pleinement que sous des conditions qui permettent un contact vivant, constant et étroit avec les étudiants. On sait combien nous en avons été loin tout au long de ces dernières années, avec des amphithéâtres-forums où le maître-gardien-de-la-science, seul devant la foule des aspirants-à-la-science, semble sorti tout droit d'une hallucinante vision de Kafka.

5. Les développements qui précèdent démontrent sans doute suffisamment la nécessité de distinguer entre le maître qui enseigne et le maître qui édifie la science, et mettent en évidence quelques uns de leurs traits caractéristiques. Il doit être bien entendu qu'il ne peut être question d'établir entre ces deux espèces de vocations une hiérarchie quelconque. Il convient simplement de réagir enfin à d'anciennes habitudes de pensée héritées du passé, qui minimisaient le rôle du premier en le subordonnant au second. Il faut donc réhabiliter la vocation pédagogique dans l'Université, en la mettant à pied d'égalité avec la vocation de recherche, et en lui fournissant les

\page{262}
conditions favorables à son épanouissement qui ont jusqu'à présent fait défaut.

En plus des maîtres eux-mêmes, les étudiants ont été évidemment les principales victimes de la survie d'un système de formation et de recrutement des maîtres basé presque exclusivement sur le critère de l'aptitude à la recherche. Il est urgent de mettre à profit dans l'immédiat la force vive libérée aujourd'hui par le mouvement de revendication étudiantes, et les prises de conscience qu'elles ont précipitées chez les enseignants, pour mettre au point des mécanismes assez souples pour distinguer et coordonner tout à la fois les deux principales vocations de l'Université, et capables de s'adapter au fur et à mesure aux besoins, changeants d'une période à l'autre et d'une Faculté à l'autre. Ces questions sont l'objet dès à présent de discussions entre étudiants et enseignants (assistants, maîtres-assistants, maîtres de conférences et professeurs), tout au moins au département de Mathématiques de la Faculté des Sciences d'Orsay, et sans doute un peu partout dans les universités en effervescence créatrice. Il faut que nous, les principaux concernés, étudiants et maîtres enseignants et chercheurs, fassions tout notre possible pour que des solutions d'ensemble soient dégagées dans le plus proche avenir, et puissent être mises en application dès la rentrée prochaine. Je passerai en revue plus bas diverses modifications en profondeur des structures et des habitudes de pensée dont la nécessité me semble dès à présent acquise. Certaines parmi les plus importantes découlent logiquement de l'état de faits que je viens d'essayer d'esquisser. Cet état de faits est bien familier sans doute à ceux de mes collègues ayant assez d'expérience de l'enseignement ou de la recherche, mais moins familier déjà, semble-t-il, aux assistants ou aux étudiants, et a fortiori aux personnes étrangères à l'Université. Mon but dans la première partie de cet exposé a été de mettre en évidence le plus clairement possible un des vices de base du système actuel, et d'inciter tous ceux qui sont concernés à réfléchir à cette question pour contribuer à l'élaboration de solutions novatrices. Les incertitudes du moment présent sont plus lourdes de promesses que toutes les certitudes sur lesquelles se reposait dame Université jusqu'à il y a peu. C'est de nous tous qui sommes concernés, étudiants, enseignants et chercheurs, non des princes qui nous gouvernent, qu'il dépend que ces promesses débouchent sur cette rénovation de l'enseignement dont personne (sujets ni princes) ne songe plus à contester la nécessité.

\page{263}
\subsection{B. Demain}

Nous esquissons ici quelques-unes des propositions qui ont été avancées, et pour la plupart discutées, concernant la structure du département de mathématique, et le système de formation, de qualification et de recrutement des enseignants et chercheurs en mathématique.

1. Le département serait partagé en deux sections, savoir l'\emph{Institut de Mathématique}, groupant les mathématiciens à vocation principale de recherche, et l'\emph{Ecole de Mathématique}, groupant les mathématiciens à vocation principale d'enseignement. Le rôle du premier serait, en plus de la recherche, l'enseignement mathématique à un niveau théorique élevé, plus particulièrement l'enseignement post-licence, destiné surtout aux futurs mathématiciens, et la formation des chercheurs. Le rôle de l'Ecole de Mathématique serait l'enseignement mathématique en général, et plus particulièrement l'enseignement pré-licence à l'usage de futurs usagers non mathématiciens. Les décisions essentielles (organisation de l'enseignement, recrutement etc.) seraient prises par une assemblée commune, l'\emph{Assemblée du département de Mathématique}, formée des maîtres de conférences et professeurs des deux sections, de représentants des étudiants et de représentants des assistants et maîtres-assistants. Entre les deux sections du département, les relations seraient très étroites, notamment par la possibilité pour un membre de l'Ecole de Mathématique de se faire détacher une année sur cinq à l'Institut de Mathématique avec dispense d'enseignement, et la possibilité pour un membre de l'Institut qui le désirerait d'être détaché, avec l'accord de ses collègues de l'Ecole de Mathématique, pendant une ou plusieurs années à cette dernière. Les membres de l'une et l'autre section seraient considérés comme exerçant des métiers essentiellement distincts, quoique étroitement liés,- ce qui n'empêcherait pas un membre de l'une des sections ayant les qualifications voulues de demander à changer de section à titre permanent, la décision incombant alors à l'Assemblée du département de Mathématique, au même titre que toute autre nomination.

\page{264}
2. Le danger de sclérose d'un enseignant, qu'il soit simple assistant ou professeur, serait dans une large mesure évité par l'usage de l'année "sabbatique" pour les membres de l'Ecole de Mathématique, mentionné plus haut, et par l'obligation faite à tout enseignant, de l'une ou l'autre section, d'un renouvellement constant de son enseignement, et notamment des sujets enseignés, tout au long de sa carrière. Il devrait être entendu notamment qu'un maître de conférences ou professeur ne pourra enseigner le même cours plus de trois années consécutives, et une \emph{commission d'arbitrage} à l'autorité incontestée devrait pouvoir être saisie lorsque certains membres de l'Assemblée du département jugeront que cette règle n'a pas été observée par un des enseignants. Enfin, pour préserver un équilibre judicieux entre les deux sections du département de Mathématique, nécessaire pour garantir la qualité de l'enseignement en même temps que le niveau scientifique du département, on pourrait prévoir une clause numérique, précisant que le nombre de maîtres de conférences et de professeurs de chaque section doit être égal au tiers au moins du nombre total des maîtres de conférences et professeurs.

3. La \emph{thèse de Mathématique} serait commune aux deux orientations "recherche" et "enseignement". Ce serait, comme par le passé, un travail écrit faisant l'objet d'un exposé oral. Il devrait présenter certaines qualités d'originalité, soit dans la mise au point ou l'exposition d'une théorie mathématique assez avancée, soit par la contribution de résultats originaux. Par le fait que des contributions originales ne sont plus exigées (*), le niveau de la thèse \uncertain{proprement} dite serait donc abaissé ; en revanche, la thèse serait un élément nécessaire mais non plus suffisant pour conférer le titre de docteur ès sciences mathématiques. Pour celui-ci, on distingue deux "mentions", la mention "recherche" et la mention "enseignement", qui qualifient respectivement pour une nomination de maître de conférences à l'Institut de Mathématique et à l'Ecole de Mathématique. Pour obtenir le doctorat mention recherche, il faudra comme par le passé, en plus d'une culture mathématique solide, avoir obtenu des résultats originaux d'une importance indiscutable. Compte tenu de l'existence d'une carrière parallèle "enseignement", il conviendrait à ce sujet d'être même plus exigeant qu'on ne l'a été souvent ces dernières années,

(*) Au lieu de supprimer entièrement cette exigence, on peut aussi essayer de se borner à l'assouplir.

\note{Le mot « proprement » est à demi effacé sur la copie ; le contexte le donne. Un caractère surfrappé avant « mention "recherche" » et deux marques en marge relèvent de la frappe.}

\page{265}
sous la pression du manque d'enseignants. Il faudrait que le docteur mention "recherche" ait une personnalité scientifique assez forte pour être jugé capable de former lui-même des chercheurs de valeur. Quant au doctorat mention enseignement, il serait décerné sur critères pédagogiques, se basant sur un minimum de quatre à cinq années de pratique de l'enseignement comme assistant et maître assistant, avec pendant une année au moins la responsabilité complète d'un cours. De plus, les activités d'enseignement du candidat devront avoir porté sur au moins trois sujets différents. Il serait tenu compte, dans l'attribution de ce doctorat, des avis des chefs de service sous la direction desquels le candidat aura été assistant, ainsi que des avis exprimés année par année par les étudiants enseignés par le candidat. Un étudiant, élu parmi ceux ayant été enseignés par le candidat, pourrait utilement faire partie du jury décidant de l'attribution du titre de docteur d'enseignement.

4. Pour que la revalorisation de la vocation d'enseignant à l'Université soit authentique, et ne se réduise à une façon pieuse de cacher une dévalorisation de la fonction de maître de conférences ou de professeur à l'Université, qu'il faut au contraire éviter à tout prix, il est essentiel que le doctorat de Mathématique (mention enseignement) ne soit décerné que sur des critères de sélection très sérieux, et qu'il ne puisse être question pour son attribution d'invoquer des arguments d'\emph{ancienneté} du candidat dans les fonctions d'assistant et de maître assistant. Ne devrait être attribué ce titre qu'à un candidat ayant les qualifications voulues pour devenir un bon maître de conférences (mention enseignement), i.e. essentiellement, un candidat ayant prouvé ses capacités pour assumer la responsabilité complète d'un cours, et donnant des garanties d'ouverture intellectuelle suffisantes pour qu'on soit fondé à penser qu'il sera en mesure d'assumer la responsabilité complète de n'importe quel cours de Mathématique d'un niveau pré-licence. Pour fonder cette conviction, en plus des appréciations des étudiants et des chefs de service sur les activités d'enseignement proprement dites, le jury pourra tenir compte de tous autres signes d'initiative intellectuelle du candidat, que ce soit dans la thèse de celui-ci, ou par sa participation active à des séminaires, la rédaction de notes polycopiées, la publication d'articles de recherche mathématiques

\page{266}
ou pédagogiques, etc. Ces mêmes éléments d'appréciation auront évidemment leur rôle à jouer pour la nomination d'un docteur d'enseignement à un poste de maître de conférences (mention enseignement), décidée (rappelons-le) par l'Assemblée de département, donc avec participation effective des délégués d'étudiants. (Pour la nomination d'un maître de conférences mention recherche, il conviendrait par contre de prévoir, en cas de demande par l'Institut de Mathématique, une cooptation par une assemblée plus restreinte, comprenant tous les membres de l'Institut de Mathématique, et où ne figuereraient pas de délégués étudiants).

5. Le point le plus délicat à l'heure actuelle semble celui du \emph{rôle dévolu à l'assistant et au maître-assistant} dans l'Université de demain. La doctrine la plus satisfaisante dans l'absolu semblerait la suivante. \emph{La fonction d'assistant et de maître-assistant serait considéré comme une fonction essentiellement} non permanente, \emph{correspondante aux dernières étapes de la formation d'un maître de conférences} (mention enseignement, le plus souvent). Elle doit donc \emph{normalement} déboucher sur un doctorat (d'enseignement, le plus souvent), et sur une nomination de maître de conférences, après une période de l'ordre de quatre à six ans. Dans cette optique, il convient donc de ne titulariser un assistant stagiaire que lorsque le département a de sérieuses raisons de penser que le candidat a les qualités requises pour aller jusqu'au doctorat, donc après la soutenance de la thèse, et sans doute pas avant deux ans de pratique de l'enseignement le plus souvent. Il faudrait donc pourvoir à la réorientation professionnelle des stagiaires assistants qui ne donneraient pas satisfaction, et, exceptionnellement, des enseignants ayant exercé les fonctions d'assistant ou maître-assistant pendant six années sans arriver à obtenir le doctorat. Ce dernier cas sera évidemment délicat du fait du statut de la fonction publique dont bénéficiera cet enseignant, mais on peut espérer qu'il sera rare. et qu'il sera possible dans la plupart des cas, avec l'accord de l'intéressé, de lui trouver un débouché en-dehors de l'enseignement, plus en rapport avec ses aptitudes réelles.

Si on compte que la durée moyenne des fonctions d'assistants ou maître-assistant avant le doctorat et la nomination à un poste de maître de conférences est de cinq années, sur une activité d'enseignement qui porte sur environ quarante ans du total (de vingt-cinq à soixante-cinq ans), on

\note{« figuereraient » : sic.}

\page{267}
trouve (dans l'optique où nous nous plaçons) que la proportion normale des assistants et maîtres-assistants dans l'ensemble des enseignants du département de Mathématique devrait être de l'ordre de un huitième. Il faut compter sans doute, compte tenu des transitions nécessaires et de la dynamique de l'accroissement régulier des besoins en effectifs d'enseignants, que pendant assez longtemps cette proportion resterait de l'ordre de un cinquième à un quart. C'est à peu près la proportion inverse de celle qui est de règle à l'heure actuelle ! Il est clair qu'un tel renversement n'est concevable que moyennant une transformation radicale des méthodes d'enseignement actuellement en vigueur à l'Université, où un petit état-major de professeurs et de maîtres de conférences dirige de haut une armée d'assistants et de maîtres-assistants dont bien peu (tout comme dans l'armée moins pacifique qui fait la fierté de la Nation) portent dans leur sac de fameux baton de Maréchal ! Une telle évolution semble liée indissolublement à la disparition du système des cours magistraux, disparition que les étudiants réclament avec insistance, mais que de nombreux professeurs auront sans doute beaucoup de mal à admettre, et même à imaginer…

6. Dans l'immédiat, en plus des obstacles psychologiques sans doute extrêmement sérieux de la part des professeurs, même parmi ceux qui sont le plus sincèrement acquis au courant de renovation universitaire parti des étudiants, il faut bien reconnaître qu'il y a également d'autres obstacles à une application conséquente de la nouvelle conception présentée ici du rôle de l'assistant et du mapitre-assistant. Il y a d'abord, bien entendu, la question des crédits, un maître-assistant "coûtant moins cher" à l'Etat qu'un maître de conférences. Les derniers évènements, il est vrai, permettent de se faire une idée si les "économies" à court terme ainsi réalisées par l'Etat sont rentables à longue échéance. Un problème sans doute plus délicat, car en grande partie indépendant du bon vouloir de l'Etat dans les années à venir, est celui de savoir combien parmi les actuels assistants et maîtres-assistants auraient un niveau intellectuel suffisant pour pouvoir, dans les deux ou trois années qui viennent, obtenir un doctorat d'enseignement et exercer les fonctions de maître de conférences (mention enseignement). Il est probable que nombre d'entre eux, (parmi ceux notamment qui s'épuisent actuellement en vains efforts pour

\note{« mapitre-assistant » : sic, faute de frappe.}

\page{268}
"décrocher" une thèse ancien style), stimulés par une option nouvelle qui s'offre à eux, porteraient à leur travail d'enseignement un intérêt plus actif et une ardeur nouvelle, dont les étudiants seraient les premiers bénéficiaires, ardeur qui de plus permettrait à ces enseignants d'atteindre à la qualification voulue pour devenir maîtres de conférences. Il serait sans doute malaisé, même à de plus compétents que moi en la matière, de faire aucun pronostic sur la proportion parmi les assistants et maîtres-assistants qui seraient en mesure de saisir l'opportunité de promotion professionnelle qui leur serait ainsi offerte. Il semble probable en tout cas que nombreux également seront ceux qui n'auront pas les moyens intellectuels nécessaires, et que pendant de longues années et jusqu'à leur retraite, ces derniers constitueront, parmi les jeunes promotions d'assistants et de maîtres-assistants faisant leur apprentissage du métier de professeur, une sorte d'arrière-garde de maîtres-assistants à vie, qui auront leur rôle à jouer dans l'inévitable transition entre l'Université d'hier et celle de demain.

7. Dans l'hypothèse où des conceptions dans l'esprit de celles esquissées ici finissent par prévaloir au cours des prochaines années (ce qui sans doute ne pourra se faire que si la contestation étudiante conserve sa vigueur, sans se reposer sur les lauriers des avantages acquis), il faudrait revoir alors des régimes de transition sans doute assez longs, caractérisés par la prédominance numérique parmi les enseignants d'un bon nombre d'assistants et de maîtres-assistants plus ou moins "immobiles", et parmi les professeurs et maîtres de conférences, de ceux qui se seront définis "chercheurs" plutôt qu'"enseignants". Parmi les nouveaux postes créés ou devenus vacants, il faudrait donc prévoir qu'une majorité devraient être attribués à des maîtres de conférences mention "enseignement", pour arriver dans un temps X à l'état d'équilibre relatif du "régime de croisière" d'un tiers du moins de maîtres de conférences mention "enseignement". Une telle politique universitaire semblerait conforme, sur le plan global, à la proportion relativement réduite des chercheurs de haut niveau parmi les jeunes mathématiciens qu'intéresse la carrière universitaire. Ellé demandera cependant, de la part des départements de Mathématique à très fort pouvoir d'attraction, tels ceux de Paris et d'Orsay, une volonté délibérée de résister à la tentation de profiter de tous les postes disponibles

\note{La copie porte ici des retouches à l'encre, de main non déterminée, signalées par des croix en marge : une virgule effacée après « assistants », une virgule ajoutée après « "immobiles" », les guillemets de « qu'"enseignants" » déplacés, et les lettres de « du » et « de » (« d'un tiers du moins de maîtres ») repassées. Le texte est donné dans son état corrigé là où la correction se lit. « Ellé » : sic.}

\page{269}
pour "thésauriser" les meilleurs chercheurs. A cette condition serait favorisé également le processus de "décentralisation des compétences", dont il a été question depuis longtemps et que chacun appelle de ses voeux, voeux plus ou moins sincères, il faut bien l'avouer, suivant que celui qui les fait se trouve en province ou à Paris. Je serais certes le dernier à sous-estimer la force d'une telle tentation pour un mathématicien passionné par sa science. Aussi j'estime que la seule façon efficace d'y remédier serait la création d'une \emph{Commission de mathématiciens interdépartements}, dont un des rôle serait de superviser les nominations, à laquelle les départements de Mathématique adhérents délégueraient une autorité véritable, y compris dans certains cas à préciser un droit de veto, dans le but de garantir une répartition raisonnable des chercheurs suivant les diverses Universités, en tenant compte à la fois de l'intérêt des divers centres de recherches effectivement existants dans les divers départements, et des désirs des chercheurs concernés. Il s'agit manifestement là d'un point extrêmement délicat, qui demandera des mécanismes d'autocontrôle d'une grande souplesse (excluant bien entendu l'intervention de toute instance ministérielle). Elles ne pourront évidemment fonctionner que moyennant une évolution considérable dans l'état d'esprit des mathématiciens concernés, et plus particulièrement des parisiens, habitués à une concentration des compétences qui depuis longtemps dépasse de très loin les besoins particuliers de chacun de nous, et qui à l'échelle nationale constitue une sorte de "gaspillage par l'accumulation".

8. Pour terminer, il convient de rappeler que jusqu'à aujourd'hui, l'autorité à l'intérieur du département a été détenue exclusivement par les maîtres de conférences et les professeurs, recrutés suivant les critères traditionnels, c'est-à-dire grosso-modo, par des chercheurs. Ces derniers constituent une véritable "aristocratie éclairée", et exercent un gouvernement absolu sur une armée d'assistants et maîtres-assistants, aux effectifs quatre ou cinq fois plus nombreux, et sur un "sous-prolétariat" d'étudiants. L'accession d'un maître-assistant à l'aristocratie régnante est soumise à des critères de recherche d'une rigueur telle que la grande majorité doivent y renoncer, et par là même renoncer à un épanouissement normal de leurs capacités pédagogiques, comme membres responsables à part entière de la communauté des ensei-

\note{« un des rôle » : sic.}

\page{270}
gnants. Quant aux étudiants, qui sont en principe la raison d'être de l'Université, leurs besoins dans tout cela ont été pratiquement ignorés. C'est là un état de faits absurde et injuste, mais consacré par le passé, et d'une inertie immense. Ceux qui détiennent toute l'autorité dans ce système, à savoir les chercheurs, n'y renonceront pour la plupart sans une résistance extrêmement sérieuse, passive ou active suivant le tempérament personnel ou les consignes partisanes, suivant la conjoncture aussi. Dans un tel état d'esprit et faute d'un vrai désir de faire du neuf, les discussions entre collègues et sans la participation d'étudiants, sur tels ou tels points d'organisation et de restructuration (comme ceux abordés dans les pages qui précèdent), risquent évidemment de devenir byzantines, et même de donner finalement aux immobilistes résolus une occasion de "noyer le poisson". Aussi il faut garder présent à l'esprit que toutes les solutions proposées sont en dernière analyse des points de détail, toujours discutables tant qu'elles ne sont pas sanctionnées et corrigées à mesure par l'expérience, dans l'esprit de la "contestation permanents" sans lequel rien qui vaille ne saurait finalement être construit. Un seul point de structure finalement est vraiment important, qui seul peut constituer une garantie relative contre la tentation des maîtres de l'Université d'enterrer au plus vite le mouvement : c'est la cogestion effective, et plus précisément, la participation des étudiants, par une représentation massive au sein des instances dirigeantes du département, à \emph{toutes} les décisions qui les concernent. Ce point acquis, et les rapports de force au niveau des décisions correspondant enfin aux rapports réels entre les fonctions des uns et des autres, enseignants et étudiants, les structures s'ajusteront d'elles-mêmes ; non en un jour ni en une année, mais chaque chose en son temps. C'est donc sur le principe et les modalités d'une cogestion effective et dans l'immédiat que tous ceux qui veulent que quelque chose change doivent, sous peine de s'éparpiller, concentrer leurs principaux efforts.

\note{Fin du tapuscrit. « contestation permanents » : sic.}

\section{Notes sur une circulaire de la commission d'Orsay}

\page{272}
\note{Pp. 272–273 : circulaire ronéotypée, collective, de la commission du département de Mathématiques de la Faculté des Sciences d'Orsay sur les deux voies d'accès au doctorat d'État (mention « recherche » et mention « enseignement »), non datée ; elle annonce la prochaine réunion « le mardi 11 juin ». Texte collectif et de tiers, non transcrit.}

\page{273}
\note{Au pied de la p. 273, sous le texte ronéotypé, des notes au crayon d'une écriture cursive rapide, que l'on attribue à Grothendieck de façon probable mais non certaine ; elles reprennent les thèmes de son propre tapuscrit (ancienneté, thèse, mandarinat). Colonne de gauche, avec deux traits horizontaux ; à droite, trois mentions entourées.}

(assistants pas à \uncertain{l'ancienneté})

option \uncertain{momentanée}

nominations par étudiants

thèse —

enseignements \uncertain{parallèles} (\uncertain{diversification})

recherche \ill{} enseignement

mandarinat

\marginal{(un peu de recherche)}

\marginal{(\uncertain{rap.} CNRS ?)}

\marginal{(décision des étudiants)}

\section{Rapport de la Commission relative au recrutement des enseignants}

\page{275}
\note{Pp. 275–280 : « Rapport de la Commission relative au recrutement des enseignants » (Faculté des Sciences d'Orsay), ronéotypie d'un manuscrit d'une autre main (pp. 275–279), la p. 280 étant une page dactylographiée « -6- » ; non daté, il renvoie à « la rentrée d'octobre 1968 ». Texte collectif et de tiers, non transcrit. En résumé : le département se partage en une École de Mathématique (premier et second cycles) et un Laboratoire de Recherche Mathématique (troisième cycle), sous une Assemblée du département comprenant assistants et représentants élus des étudiants ; le doctorat mention « enseignement » ou « recherche » qualifie pour l'une ou l'autre section, les nominations étant décidées par l'Assemblée ou des commissions qu'elle désigne, avec un droit de veto étudiant pour l'École ; des détachements périodiques et des changements de section sont prévus entre les deux.}

\end{document}
